\section{Introduction}
\label{sec:intro}

Empirical study of board game design has been hampered by the absence of a large-scale,
consistently-scored corpus: existing evaluations are either narrow in scope, inconsistent
in methodology, or rely on community aggregate ratings that conflate popularity with quality.
This paper addresses that gap by introducing the TIGRIS corpus --- 150 games scored on 32
axes by independent evaluators under a structured protocol --- and establishing the
reliability infrastructure required to use the corpus as an empirical foundation for
subsequent analysis.

\subsection{Motivation}

Board games represent a rich but under-studied domain for design research: they are
discrete, rule-governed artifacts with well-defined mechanical components, yet they produce
emergent social and strategic experiences that resist simple decomposition.
Prior quantitative work has largely relied on BoardGameGeek's community-assigned weight
scores~\cite{woods2012eurogames}, which capture perceived complexity but not design
character, or on single-axis expert rankings that cannot support multivariate analysis.
The TIGRIS corpus is designed to fill this gap: 32 axes covering tension, interaction
architecture, elegance, range, and accessibility dimensions, scored at a granularity
sufficient to distinguish games within the same weight band.

\subsection{Contributions}

This paper makes three contributions: (1) the TIGRIS corpus itself, released as an open
dataset; (2) a two-mode scoring protocol distinguishing fast-track estimates from
full-pipeline scores, with explicit reliability thresholds for each mode; and (3)
per-axis Cronbach's $\alpha$ reliability reports establishing the psychometric foundation
for downstream analysis.
